\SourcePage{78}{81}
\section{Flux Density}

In classical mechanics a particle's velocity $\mathbf v$ is related to its
momentum by $\mathbf p=m\mathbf v$. As expected, quantum mechanics has the
same relation between the corresponding operators. This follows at once by
calculating $\hat{\mathbf v}=\hat{\dot{\mathbf r}}$ from the general rule
(9.2) for differentiating operators with respect to time:
\[
  \hat{\mathbf v}=\frac{i}{\hbar}
  \left(\hat H\mathbf r-\mathbf r\hat H\right).
\]
Using (17.5) for $\hat H$ together with (16.5),%
\SourcePage{79}{82}
we obtain
\begin{equation}
  \hat{\mathbf v}=\hat{\mathbf p}/m.
  \tag{19.1}
\end{equation}
The same relations plainly hold between the eigenvalues of velocity and
momentum, and between their mean values in any state.

Like momentum, the velocity of a particle cannot have a definite value
simultaneously with its coordinates. Velocity multiplied by an infinitesimal
time element $dt$, however, determines the particle's displacement during
$dt$. Thus the nonexistence of velocity simultaneously with coordinates means
that if a particle occupies a definite point of space at one instant, it no
longer has a definite position at the next, infinitesimally close instant.

We record a useful formula for the operator $\hat{\dot f}$, the time
derivative of a quantity $f(\mathbf r)$ which is a function of the particle's
radius vector. Since $f$ commutes with $U(\mathbf r)$, we find
\[
  \hat{\dot f}=\frac{i}{\hbar}(\hat Hf-f\hat H)
  =\frac{i}{2m\hbar}(\hat{\mathbf p}^{2}f-f\hat{\mathbf p}^{2}),
\]
while (16.4) gives
\[
  \hat{\mathbf p}^{2}f
  =\hat{\mathbf p}\left(f\hat{\mathbf p}-i\hbar\Grad{f}\right),
  \qquad
  f\hat{\mathbf p}^{2}
  =\left(\hat{\mathbf p}f+i\hbar\Grad{f}\right)\hat{\mathbf p}.
\]
Hence the desired expression is
\begin{equation}
  \hat{\dot f}=\frac{1}{2m}
  \left(\hat{\mathbf p}\,\Grad{f}+\Grad{f}\cdot\hat{\mathbf p}\right).
  \tag{19.2}
\end{equation}

Next, let us find the acceleration operator. We have
\[
  \hat{\dot{\mathbf v}}
  =\frac{i}{\hbar}(\hat H\hat{\mathbf v}-\hat{\mathbf v}\hat H)
  =\frac{i}{m\hbar}(\hat H\hat{\mathbf p}-\hat{\mathbf p}\hat H)
  =\frac{i}{m\hbar}(U\hat{\mathbf p}-\hat{\mathbf p}U).
\]
Using (16.4), we find
\begin{equation}
  m\hat{\dot{\mathbf v}}=-\Grad{U}.
  \tag{19.3}
\end{equation}
In form, this operator equation coincides exactly with the equation of motion
(Newton's equation) of classical mechanics.

The integral $\int_V|\Psi|^2\,dV$ over a finite volume $V$ is the probability
of finding the particle in that volume. Let us calculate its time derivative:
\[
  \frac{d}{dt}\int_V|\Psi|^2\,dV
  =\int_V\left(\Psi\frac{\partial\Psi^*}{\partial t}
  +\Psi^*\frac{\partial\Psi}{\partial t}\right)dV
  =\frac{i}{\hbar}\int_V
  \left(\Psi\hat H^*\Psi^*-\Psi^*\hat H\Psi\right)dV.
\]
\SourcePage{80}{83}
Substitution of
\[
  \hat H=\hat H^*=-\frac{\hbar^2}{2m}\Delta+U(x,y,z)
\]
and use of the identity
\[
  \Psi\Delta\Psi^*-\Psi^*\Delta\Psi
  =\Div{(\Psi\Grad{\Psi^*}-\Psi^*\Grad{\Psi})}
\]
give
\[
  \frac{d}{dt}\int_V|\Psi|^2\,dV
  =-\int_V\Div{\mathbf j}\,dV,
\]
where $\mathbf j$ denotes the vector\footnote{If $\psi$ is written as
$|\psi|e^{i\alpha}$, then
\begin{equation}
  \mathbf j=\frac{\hbar}{m}|\psi|^2\Grad{\alpha}.
  \tag{19.4a}
\end{equation}}
\begin{equation}
  \mathbf j=\frac{i\hbar}{2m}
  (\Psi\Grad{\Psi^*}-\Psi^*\Grad{\Psi})
  =\frac{1}{2m}
  (\Psi\hat{\mathbf p}^{*}\Psi^*+\Psi^*\hat{\mathbf p}\Psi).
  \tag{19.4}
\end{equation}
By Gauss's theorem, the integral of $\Div{\mathbf j}$ can be transformed into
an integral over the closed surface bounding $V$:
\begin{equation}
  \frac{d}{dt}\int_V|\Psi|^2\,dV
  =-\oint\mathbf j\,d\mathbf f.
  \tag{19.5}
\end{equation}
It follows that $\mathbf j$ may be called the \emph{probability-flux-density
vector}, or simply the \emph{flux density}. Its surface integral is the
probability that the particle crosses that surface per unit time. The vector
$\mathbf j$ and the probability density $|\Psi|^2$ obey
\begin{equation}
  \frac{\partial|\Psi|^2}{\partial t}+\Div{\mathbf j}=0,
  \tag{19.6}
\end{equation}
an equation analogous to the classical continuity equation.

The free-motion wave function---the plane wave (17.9)---can be normalized so
as to describe a particle flux of unit density (a flux in which, on average,
one particle per unit time passes through a unit area of its cross-section).
The function is
\begin{equation}
  \Psi=\frac{1}{\sqrt v}
  \exp\left[-\frac{i}{\hbar}(Et-\mathbf p\cdot\mathbf r)\right],
  \tag{19.7}
\end{equation}
where $v$ is the particle's speed. Indeed, substituting this function into
(19.4) gives $\mathbf j=\mathbf p/mv$, the unit vector in the direction of
motion.
\SourcePage{81}{84}
It is useful to show how the mutual orthogonality of wave functions of states
with different energies follows directly from the Schrödinger equation. Let
$\psi_m$ and $\psi_n$ be two such functions; they satisfy
\[
  -\frac{\hbar^2}{2m}\Delta\psi_m+U\psi_m=E_m\psi_m,
\]
\[
  -\frac{\hbar^2}{2m}\Delta\psi_n^*+U\psi_n^*=E_n\psi_n^*.
\]
Multiply the first equation by $\psi_n^*$ and the second by $\psi_m$, and
subtract them term by term. This gives
\[
  (E_m-E_n)\psi_m\psi_n^*
  =\frac{\hbar^2}{2m}(\psi_m\Delta\psi_n^*-\psi_n^*\Delta\psi_m)
  =\frac{\hbar^2}{2m}\Div{
  (\psi_m\Grad{\psi_n^*}-\psi_n^*\Grad{\psi_m})}.
\]
Integrating both sides over all space, the right-hand side vanishes after
application of Gauss's theorem, and we obtain
\[
  (E_m-E_n)\int\psi_m\psi_n^*\,dV=0.
\]
Since $E_m\ne E_n$ by assumption, the required orthogonality relation follows:
\[
  \int\psi_m\psi_n^*\,dV=0.
\]
